%% ===========================================================================
%% SCAFF: the stage-crossover attribution layer.
%%
%% Inserted into Section~\ref{sec:deltas} immediately AFTER \subsubsection{Neutral and
%% Sensitive Trajectories} and BEFORE \subsubsection{Invariance and Quality Benefits}, so that
%% the per-stage deltas that follow read as choices of the distance function $d_s$ rather than
%% as the framework's primitive.
%%
%% Requires in the preamble:
%%   \newcommand{\Dnat}{D^{\mathrm{nat}}}
%%   \newcommand{\DA}{D^{A}}
%%   \newcommand{\DZ}{D^{Z}}
%% ===========================================================================

\subsubsection{Stage Crossover: Separating Direct Sensitivity from Inherited Dependence}
\label{sec:scaff}

Comparing $O_s(\tau^{a}_c)$ against $O_s(\tau_c)$ at every stage localizes \emph{where a
trajectory diverges}. It does not localize \emph{where the protected attribute entered}, and in a
stateful pipeline those are different questions. Each stage consumes the output of its
predecessor, so a stage that never inspects the attribute still produces a different output once
the state it inherits has been changed upstream. A stage-wise comparison of outputs therefore
reports every stage downstream of a violation as attribute-sensitive, and the resulting diagnosis
implicates the stages that merely transmitted the difference alongside the one that created it.
The distinction is not academic: it determines which component an operator should repair.

We resolve this by intervening on the two inputs of a stage separately. Fix a stage $s$ and write
$Z_s$ for its \emph{legitimate parent state}---the complete non-protected input it receives from
upstream, which for the substrate of Section~\ref{sec:substrate} is the accumulated preference
state at \stage{Retrieve}, that state together with the candidate list at \stage{Rank}, and so on.
By construction $Z_s$ never contains the attribute. Let $Z^{a}_s$ and $Z^{a'}_s$ denote the parent
states that the two protected conditions \emph{naturally} reach, and let $f_s$ be the stage's
transition with $\phi_s$ its audited projection. Define the crossover outcome
\begin{equation}
  Y^{i,j}_s = \phi_s\!\left(f_s(Z^{i}_s,\, A = j;\, \epsilon)\right),
  \qquad i, j \in \{a, a'\},
  \label{eq:crossover}
\end{equation}
in which the first index selects \emph{whose inherited state} the stage is supplied and the second
selects \emph{which descriptor the stage itself is shown}. The four cells of
Equation~\eqref{eq:crossover} are the $2 \times 2$ design of this framework: the two diagonal
cells are the naturally evolving trajectories, and the two off-diagonal cells are counterfactuals
that no end-to-end replay produces, because they require a stage to be re-executed on an input its
own arm never reached.

Three quantities follow, all expressed through the same stage-specific distance $d_s$. Natural
divergence,
\begin{equation}
  \Dnat_s(c, a) = d_s\!\left(Y^{a,a}_s,\, Y^{a',a'}_s\right),
  \label{eq:dnat}
\end{equation}
is the diagonal contrast, and is exactly what a stage-wise comparison of paired trajectories
measures. Direct sensitivity holds the parent state fixed and swaps only the descriptor,
\begin{equation}
  \DA_s(c, a) = \tfrac{1}{2}\left[
    d_s\!\left(Y^{a,a}_s,\, Y^{a,a'}_s\right) +
    d_s\!\left(Y^{a',a}_s,\, Y^{a',a'}_s\right)\right],
  \label{eq:ddirect}
\end{equation}
averaging over the two parent states so that the estimate does not depend on which arm is treated
as the reference. Inherited dependence holds the descriptor fixed and swaps the parent state,
\begin{equation}
  \DZ_s(c, a) = \tfrac{1}{2}\left[
    d_s\!\left(Y^{a,a}_s,\, Y^{a',a}_s\right) +
    d_s\!\left(Y^{a,a'}_s,\, Y^{a',a'}_s\right)\right].
  \label{eq:dinherited}
\end{equation}

The interpretation of each is fixed by what it holds constant. A non-zero $\DA_s$ can only arise
from the stage's own response to the descriptor, since nothing else about its input changed; this
is the quantity that attributes a violation to a location. A non-zero $\DZ_s$ arises when the
stage responds to a difference in what it was given, which is what a correctly functioning stage
is supposed to do---the disparity is real and user-visible, but the stage is not where it should be
corrected. $\Dnat_s$ confounds the two, and Equations~\eqref{eq:ddirect}
and~\eqref{eq:dinherited} are how it is decomposed. Under a deterministic stage the three are
related by the triangle inequality rather than additively, so we report them as three coordinates
and do not combine them into a scalar.

We also report a fourth coordinate, the within-condition distance obtained by repeating a single
arm under identical inputs, which estimates the noise floor $\epsilon$ contributes. This is not
decoration: a stage whose $\DA_s$ is smaller than its own repeat variance has not been shown to be
attribute-sensitive at all, and a decision rule applied without it converts decoding
stochasticity into findings.

\paragraph{Reading the coordinates}
Fixing a margin $\tau$ above the noise floor, each stage is assigned one of four diagnoses:
\emph{direct} when $\DA_s > \tau$; \emph{inherited} when $\DA_s \leq \tau$ while $\Dnat_s$ and
$\DZ_s$ both exceed $\tau$; \emph{mixed} when $\DA_s$ and $\DZ_s$ both exceed $\tau$, the case in
which a stage both reads the attribute and receives an already-corrupted state; and
\emph{invariant} otherwise. The diagnosis is defined per conversation and attribute value. A
corpus mean over these coordinates is a summary of a distribution and is not itself a verdict:
a stage that is directly sensitive on a minority of conversations will average below any useful
margin, which is why Section~\ref{sec:scaff-results} reports the distribution of per-conversation
diagnoses alongside the means.

\paragraph{Three claims, separated}
Equations~\eqref{eq:ddirect} and~\eqref{eq:dinherited} let us keep apart three assertions that
the fairness literature routinely merges, and that carry different evidential burdens.
\emph{Direct sensitivity} is an empirical claim about a stage's transition function, established
by $\DA_s$ and requiring no ground truth. \emph{Procedural unfairness} is the further claim that
this particular dependence is prohibited for this particular stage---a normative premise supplied
by the auditor, not measured. It is stated per stage because it varies per stage: a descriptor may
be a legitimate input to a safety filter and an illegitimate one to preference inference.
\emph{Disparate harm} is the strongest claim, that the affected party is worse off, and it alone
requires a reference $\Omega_s$ against which one output can be called better than the other:
whether an additional clarification was informative, whether a candidate set had lower recall,
whether a ranking had lower $\NDCG@K$, whether a retained memory fact was supported.

This ordering has a consequence for what an audit may report. Where a defensible $\Omega_s$
exists, all three claims can be made. Where none does, sensitivity and procedural status remain
reportable and harm must be recorded as \emph{not identified}. We take that as the correct output
rather than a gap to be filled by an LLM judge: a judge asked to adjudicate whether one clarifying
question was more useful than another, absent any reference, produces a number with no
truth-conditions attached. The judge protocol of Section~\ref{sec:judge} is accordingly scoped to
semantic constructs for which agreement rather than correctness is the property of interest---
whether two clarifying questions target the same facet, whether an explanation asserts something
the trace does not support, whether framing is stereotypical---and deterministic checks and task
ground truth take precedence wherever both are available.

\paragraph{Relation to the per-stage deltas}
The deltas of Section~\ref{sec:stage-deltas} are not superseded. Each supplies the distance $d_s$
for its stage---clarification-demand difference at \stage{Elicit}, Jaccard divergence of the
candidate pool at \stage{Retrieve}, a ranking contrast at \stage{Rank}, coverage and leakage at
\stage{Explain}, profile divergence at \stage{Memory}---and each is evaluated at all four cells of
Equation~\eqref{eq:crossover} rather than at the diagonal alone. The framework's earlier
formulation is thus the special case in which only $\Dnat_s$ is computed, and the group statistics
of Section~\ref{sec:groupmetrics} apply unchanged to whichever coordinate is being aggregated,
provided the coordinate is named: $\SNSR_s$ computed on $\Dnat_s$ and on $\DA_s$ answer different
questions, and reporting one under the other's name is the specific error this section exists to
prevent.

\paragraph{Repair as a test of the attribution}
An attribution claim should predict the effect of acting on it, and this one does. Replace the
implicated stage's output on one arm with the corresponding output from the other, re-run every
descendant stage on the patched state, and measure the resulting change in the final endpoint
divergence. We call the reduction \emph{removability}. If $\DA_s$ has identified the origin, then
repairing at $s$ should remove the downstream disparity while repairing at a stage the crossover
exonerates should not, and the gap between those two numbers is a direct test of the decomposition
rather than a restatement of it. One constraint on the experiment is worth stating because
violating it produces a spurious success: removability must not be measured on the output of the
stage being repaired, since substituting that output sets the measured divergence to zero
irrespective of where the fault originated. Removability is therefore evaluated strictly
downstream of the repaired stage, which for a repair at \stage{Rank} in a single-turn setting
means at the following turn.

\paragraph{Why endpoint auditing cannot substitute}
Two consequences of the decomposition explain why an output-only audit is not a cheaper
approximation of this one. First, empirically, a trajectory can be process-sensitive while its
final ranked list stays within any reasonable equivalence margin, because the effect of an
upstream difference need not survive ranking and truncation; those cases are recorded as fair by
an endpoint audit. Second, and structurally, faults at \stage{Explain} and \stage{Memory} lie
downstream of the ranked list and therefore \emph{cannot} change it within a turn at all. For
those two stages an endpoint audit is not underpowered but blind by construction, and the
persistence of a \stage{Memory} fault only becomes observable on a subsequent turn---which is why
the multi-turn substrate of Section~\ref{sec:substrate} is load-bearing rather than incidental.
